Los anexos constituyen otro capítulo. Con base en el Diccionario de la Lengua Española, se entiende que es un agregado, es decir, aquello importante para entender el contexto del trabajo. Es necesario redactar un párrafo de introducción que describa el contenido de los anexos. No es necesario explicarlos, solo mencionar qué se incluye en este capítulo, por ejemplo:  

\begin{itemize}
    \item[] En estos anexos, se incluye información adicional para entender el contexto del trabajo. La primera parte presenta [...]. Luego, se muestran [...]. Además, se incluye [...]. Por último, se proporciona [...]-

\end{itemize}

Esto guiará al lector, ya que no debe asumirse que sabe lo mismo que el autor, por algo está leyendo el trabajo. Así que sólo copiar y pegar una serie de figuras y cuadros sin un contexto es una falta de cortesía.

Los anexos incluyen material relevante para mayor claridad y profundidad de la investigación, pero distraen al lector si se colocan en el texto del informe. Pueden agregarse cartas, manuales, guías, etc. Es necesario asignar un número a cada anexo de acuerdo con el orden de mención en el desarrollo del trabajo. 

\section{Sección de ejemplo}

Esta corresponde a una sección de ejemplo dentro de los Anexos.